% GENERATED FILE. Edit canonical lessons, then run npm run generate:books.
% canonical-source-hash: fnv1a64:d9804bda97a314ea
% canonical-lessons: GE-C504-umgraben, GE-C504-spaten, GE-C504-saeen, GE-C504-samen, GE-C504-ernte

\chapter{Säen und umgraben}
\label{ch:ge-saen-und-umgraben}
\hlchaptermodality{504}

\begin{chapteropening}
\textbf{By the end of this chapter:} \emph{I can talk about growing things.}

Complete the last lesson of chapter 504: I can talk about growing things.
\end{chapteropening}

\section[umgraben]{umgraben — to dig over}

\label{lesson:GE-C504-umgraben}



\begin{quote}
Before the new one: say the German for watering can, then the German for compost.
\end{quote}

\subsection*{You'll want to know: umgraben}
\textbf{umgraben} — \textquotedblleft{}to dig over\textquotedblright{}.

\begin{quote}
Im Frühling graben wir das Beet um.
\end{quote}

\begin{quote}
\emph{In spring we dig over the flower bed.}
\end{quote}

Separable: \emph{er gräbt um}, \emph{hat umgegraben}.

\begin{cousinweb}
\emph{umgraben} is built on \emph{um} \textquotedblleft{}around, over\textquotedblright{} and \emph{graben} \textquotedblleft{}to dig\textquotedblright{}, inherited from Old High German \emph{graban}, from Proto-Germanic \emph{grabaną} (EtymDB).
\end{cousinweb}

\begin{grammarlens}[title={What you've built}]
One more word for telling, explaining and describing at B1.
\end{grammarlens}

\subsection*{Your turn}
\begin{itemize}
\raggedright
  \item \emph{Say it:} \emph{umgraben}
  \item \emph{Say it:} \emph{umgraben}, once more
  \item \emph{Say it:} a sentence of your own with \emph{umgraben}
  \item \emph{Recall:} say the German for watering can, then the German for compost, then say \emph{umgraben} again
\end{itemize}

\begin{tcolorbox}[breakable,colback=teal!4,colframe=teal!35!black,title={Before you move on}]
What does umgraben mean? (\textquotedblleft{}To dig over\textquotedblright{}.) Say it once more.
\end{tcolorbox}
\section[der Spaten]{der Spaten — spade}

\label{lesson:GE-C504-spaten}



\begin{quote}
Before the new one: say the German for compost, then the German for to dig over.
\end{quote}

\subsection*{You'll want to know: der Spaten}
\textbf{der Spaten} — \textquotedblleft{}spade\textquotedblright{}.

\begin{quote}
Mit dem Spaten gräbt er ein Loch für den Baum.
\end{quote}

\begin{quote}
\emph{He digs a hole for the tree with the spade.}
\end{quote}

Plural: \emph{die Spaten}.

\begin{cousinweb}
\emph{Spaten} is usually traced to Low German \emph{spade}; English \emph{spade} is a cousin. EtymDB lists the German word without a source.
\end{cousinweb}

\begin{grammarlens}[title={What you've built}]
One more word for telling, explaining and describing at B1.
\end{grammarlens}

\subsection*{Your turn}
\begin{itemize}
\raggedright
  \item \emph{Say it:} \emph{der Spaten}
  \item \emph{Say it:} \emph{der Spaten}, once more
  \item \emph{Say it:} a sentence of your own with \emph{der Spaten}
  \item \emph{Recall:} say the German for compost, then the German for to dig over, then say \emph{der Spaten} again
\end{itemize}

\begin{tcolorbox}[breakable,colback=teal!4,colframe=teal!35!black,title={Before you move on}]
What does der Spaten mean? (\textquotedblleft{}Spade\textquotedblright{}.) Say it once more.
\end{tcolorbox}
\section[säen]{säen — to sow}

\label{lesson:GE-C504-saeen}



\begin{quote}
Before the new one: say the German for to dig over, then the German for spade.
\end{quote}

\subsection*{You'll want to know: säen}
\textbf{säen} — \textquotedblleft{}to sow\textquotedblright{}.

\begin{quote}
Im April säen wir Karotten und Salat.
\end{quote}

\begin{quote}
\emph{In April we sow carrots and lettuce.}
\end{quote}

Regular: \emph{hat gesät}.

\begin{cousinweb}
\emph{säen} is inherited from Old High German \emph{sāen}, from Proto-Germanic \emph{sēaną} \textquotedblleft{}to sow\textquotedblright{} (EtymDB); English \emph{sow} is a cousin.
\end{cousinweb}

\begin{grammarlens}[title={What you've built}]
One more word for telling, explaining and describing at B1.
\end{grammarlens}

\subsection*{Your turn}
\begin{itemize}
\raggedright
  \item \emph{Say it:} \emph{säen}
  \item \emph{Say it:} \emph{säen}, once more
  \item \emph{Say it:} a sentence of your own with \emph{säen}
  \item \emph{Recall:} say the German for to dig over, then the German for spade, then say \emph{säen} again
\end{itemize}

\begin{tcolorbox}[breakable,colback=teal!4,colframe=teal!35!black,title={Before you move on}]
What does säen mean? (\textquotedblleft{}To sow\textquotedblright{}.) Say it once more.
\end{tcolorbox}
\section[der Samen]{der Samen — seed}

\label{lesson:GE-C504-samen}



\begin{quote}
Before the new one: say the German for spade, then the German for to sow.
\end{quote}

\subsection*{You'll want to know: der Samen}
\textbf{der Samen} — \textquotedblleft{}seed\textquotedblright{}.

\begin{quote}
Die Samen brauchen viel Wasser und Sonne.
\end{quote}

\begin{quote}
\emph{The seeds need lots of water and sun.}
\end{quote}

Plural: \emph{die Samen}.

\begin{cousinweb}
\emph{Samen} is inherited from Middle High German \emph{sām} (EtymDB), built on the same Germanic root as \emph{säen} \textquotedblleft{}to sow\textquotedblright{}.
\end{cousinweb}

\begin{grammarlens}[title={What you've built}]
One more word for telling, explaining and describing at B1.
\end{grammarlens}

\subsection*{Your turn}
\begin{itemize}
\raggedright
  \item \emph{Say it:} \emph{der Samen}
  \item \emph{Say it:} \emph{der Samen}, once more
  \item \emph{Say it:} a sentence of your own with \emph{der Samen}
  \item \emph{Recall:} say the German for spade, then the German for to sow, then say \emph{der Samen} again
\end{itemize}

\begin{tcolorbox}[breakable,colback=teal!4,colframe=teal!35!black,title={Before you move on}]
What does der Samen mean? (\textquotedblleft{}Seed\textquotedblright{}.) Say it once more.
\end{tcolorbox}
\section[die Ernte]{die Ernte — harvest}

\label{lesson:GE-C504-ernte}



\begin{quote}
Before the new one: say the German for to sow, then the German for seed.
\end{quote}

\subsection*{You'll want to know: die Ernte}
\textbf{die Ernte} — \textquotedblleft{}harvest\textquotedblright{}.

\begin{quote}
Die Ernte war dieses Jahr sehr gut.
\end{quote}

\begin{quote}
\emph{The harvest was very good this year.}
\end{quote}

Plural: \emph{die Ernten}. The verb: \emph{ernten} — to harvest.

\begin{cousinweb}
\emph{Ernte} is usually traced to Old High German \emph{arnōt} \textquotedblleft{}harvest\textquotedblright{}, built on an old Germanic word for the harvest season; EtymDB lists the German word without a source.
\end{cousinweb}

\begin{grammarlens}[title={What you've built}]
One more word for telling, explaining and describing at B1.
\end{grammarlens}

\subsection*{Your turn}
\begin{itemize}
\raggedright
  \item \emph{Say it:} \emph{die Ernte}
  \item \emph{Say it:} \emph{die Ernte}, once more
  \item \emph{Say it:} a sentence of your own with \emph{die Ernte}
  \item \emph{Recall:} say the German for to sow, then the German for seed, then say \emph{die Ernte} again
\end{itemize}

\begin{tcolorbox}[breakable,colback=teal!4,colframe=teal!35!black,title={Before you move on}]
What does die Ernte mean? (\textquotedblleft{}Harvest\textquotedblright{}.) Say it once more.
\end{tcolorbox}
